\usemodule[memos][fontfamily=adobehans,layout=narrow]
\setupinteractionscreen[option=bookmark]

\usemodule[mindmap]


\defineframed[RightCodeExampleFramedX][
    frame=off,before=,after=,
    minheight=2\baselineskip,
    align=top,
    offset=none,frameoffset=0pt,
    width=0.5\textwidth,
    loffset=.5em]

\definebuffer[MindmapExample]

\permanent\protected\def\stopMindmapExample{%
    \noindent\placesidebyside
    {\LeftCodeExampleFramedX{\typeMindmapExample\vfil}}
    {\RightCodeExampleFramedX{\bgroup\getMindmapExample\egroup\vfil}}}

\dontcomplain
\definecolor[lightblue][r=0.68, g=0.85, b=0.90]
\definecolor[lightgreen][r=0.56, g=0.93, b=0.56]
\definecolor[lightyellow][r=1.00, g=1.00, b=0.88]
\definecolor[lightpink][r=1.00, g=0.71, b=0.76]
\definecolor[lightcyan][r=0.88, g=1.00, b=1.00]
\definecolor[lightgray][r=0.83, g=0.83, b=0.83]
\definecolor[lightlavender][r=0.90, g=0.90, b=1.00]
\definecolor[darkblue][r=0.00, g=0.00, b=0.55]
\definecolor[darkgreen][r=0.00, g=0.39, b=0.00]
\definecolor[darkorange][r=1.00, g=0.55, b=0.00]
\definecolor[darkred][r=0.55, g=0.00, b=0.00]
\definecolor[darkcyan][r=0.00, g=0.55, b=0.55]
\definecolor[darkgray][r=0.33, g=0.33, b=0.33]
\definecolor[gray][r=0.50, g=0.50, b=0.50]
\definecolor[transparent][t=1]
\placebookmarks [subject,subsubject][force=yes]

\starttext

\title{思维导图模块 - JSON 语法示例}

\subject{一、基础示例}

\subsubject{1.1 最简单的思维导图}

\startMindmapExample
\mmjson{
    direction: "top-down";
    name: "中心主题";
    children: [
        {name: "子主题1"},
        {name: "子主题2"},
        {name: "子主题3"}
    ]
}
\stopMindmapExample

\subsubject{1.2 带颜色的节点}

\startMindmapExample
\mmjson{
    direction: "top-down";
    name: "项目计划";
    background_color: "lightblue";
    border_color: "darkblue";
    children: [
        {name: "需求分析"; background_color: "lightgreen"; border_color: "darkgreen"},
        {name: "设计阶段"; background_color: "lightyellow"; border_color: "darkorange"},
        {name: "开发实现"; background_color: "lightpink"; border_color: "darkred"},
        {name: "测试部署"; background_color: "lightcyan"; border_color: "darkcyan"}
    ]
}
\stopMindmapExample

\subsubject{1.3 多层结构}

\startMindmapExample
\mmjson{
    direction: "top-down";
    name: "知识体系";
    background_color: "lightblue";
    children: [
        {
            name: "编程语言";
            background_color: "lightgreen";
            children: [
                {name: "Python"; background_color: "white"},
                {name: "JavaScript"; background_color: "white"},
                {name: "PHP"; background_color: "white"}
            ]
        },
        {
            name: "算法数据结构";
            background_color: "lightyellow";
            children: [
                {name: "排序算法"; background_color: "white"},
                {name: "树结构"; background_color: "white"}
            ]
        },
        {
            name: "系统设计";
            background_color: "lightpink";
            children: [
                {name: "架构设计"; background_color: "white"},
                {name: "数据库设计"; background_color: "white"}
            ]
        }
    ]
}
\stopMindmapExample

\subject{二、布局方向}

\subsubject{2.1 四种布局方向}

\startMindmapExample
\mmjson{direction: "top-down"; name: "从上到下"; children: [{name: "子1"}, {name: "子2"}]}
\stopMindmapExample

\startMindmapExample
\mmjson{direction: "bottom-up"; name: "从下到上"; children: [{name: "子1"}, {name: "子2"}]}
\stopMindmapExample

\startMindmapExample
\mmjson{direction: "left-right"; name: "从左到右"; children: [{name: "子1"}, {name: "子2"}]}
\stopMindmapExample

\startMindmapExample
\mmjson{direction: "right-left"; name: "从右到左"; children: [{name: "子1"}, {name: "子2"}]}
\stopMindmapExample

\subsubject{2.2 径向布局}

\startMindmapExample
\mmjson{
    layout: "radial";
    radial_spacing: 60;
    name: "技术栈";
    background_color: "lightblue";
    children: [
        {name: "前端"; background_color: "lightgreen"},
        {name: "后端"; background_color: "lightyellow"},
        {name: "数据库"; background_color: "lightpink"},
        {name: "DevOps"; background_color: "lightcyan"},
        {name: "测试"; background_color: "lightlavender"}
    ]
}
\stopMindmapExample

\subject{三、节点样式}

\subsubject{3.1 形状 (shape)}

\startMindmapExample
\mmjson{
    direction: "top-down";
    name: "矩形";
    shape: "rectangle";
    background_color: "lightblue";
    children: [
        {name: "圆形"; shape: "circle"; background_color: "lightgreen"},
        {name: "椭圆"; shape: "ellipse"; background_color: "lightyellow"},
        {name: "菱形"; shape: "diamond"; background_color: "lightpink"},
        {name: "六边形"; shape: "hexagon"; background_color: "lightcyan"}
    ]
}
\stopMindmapExample

\subsubject{3.2 圆角矩形 (rounded + corner_radius)}

\startMindmapExample
\mmjson{
    direction: "top-down";
    name: "圆角矩形";
    shape: "rounded";
    corner_radius: 10;
    background_color: "lightblue";
    children: [
        {name: "小圆角"; shape: "rounded"; corner_radius: 5; background_color: "lightgreen"},
        {name: "中圆角"; shape: "rounded"; corner_radius: 10; background_color: "lightyellow"},
        {name: "大圆角"; shape: "rounded"; corner_radius: 20; background_color: "lightpink"}
    ]
}
\stopMindmapExample

\subsubject{3.3 字体大小和颜色 (font_size, text_color)}

\startMindmapExample
\mmjson{
    direction: "top-down";
    name: "大标题";
    font_size: 14;
    text_color: "darkblue";
    background_color: "lightblue";
    children: [
        {name: "中标题"; font_size: 12; text_color: "darkgreen"; background_color: "lightgreen"},
        {name: "小标题"; font_size: 10; text_color: "darkorange"; background_color: "lightyellow"},
        {name: "更小"; font_size: 8; text_color: "darkred"; background_color: "lightpink"}
    ]
}
\stopMindmapExample

\subsubject{3.4 文本对齐 (text_align)}

\startMindmapExample
\mmjson{
    direction: "top-down";
    name: "左对齐";
    text_align: "left";
    width: 100;
    background_color: "lightblue";
    children: [
        {name: "居中"; text_align: "center"; width: 100; background_color: "lightgreen"},
        {name: "右对齐"; text_align: "right"; width: 100; background_color: "lightyellow"}
    ]
}
\stopMindmapExample

\subsubject{3.5 边框样式 (border_width, border_color, border_direction)}

\startMindmapExample
\mmjson{
    direction: "top-down";
    name: "粗边框";
    border_width: 3;
    border_color: "darkblue";
    background_color: "lightblue";
    children: [
        {name: "中边框"; border_width: 2; border_color: "darkgreen"; background_color: "lightgreen"},
        {name: "细边框"; border_width: 1; border_color: "darkorange"; background_color: "lightyellow"},
        {name: "无边框"; border_width: 0; background_color: "lightpink"}
    ]
}
\stopMindmapExample

\startMindmapExample
\mmjson{
    direction: "top-down";
    name: "完整边框";
    border_direction: "all";
    background_color: "lightblue";
    border_color: "darkblue";
    children: [
        {name: "上边框"; border_direction: "top"; background_color: "lightgreen"; border_color: "darkgreen"},
        {name: "下边框"; border_direction: "bottom"; background_color: "lightyellow"; border_color: "darkorange"},
        {name: "左右边框"; border_direction: "left, right"; background_color: "lightpink"; border_color: "darkred"}
    ]
}
\stopMindmapExample

\subject{四、连接线样式}

\subsubject{4.1 连接类型 (connection)}

\startMindmapExample
\mmjson{
    direction: "top-down";
    name: "直线连接";
    connection: "straight";
    background_color: "lightblue";
    children: [
        {name: "子节点1"; background_color: "lightgreen"},
        {name: "子节点2"; background_color: "lightyellow"},
        {name: "子节点3"; background_color: "lightpink"}
    ]
}
\stopMindmapExample

\startMindmapExample
\mmjson{
    direction: "top-down";
    name: "曲线连接";
    connection: "curved";
    background_color: "lightblue";
    children: [
        {name: "子节点1"; background_color: "lightgreen"},
        {name: "子节点2"; background_color: "lightyellow"},
        {name: "子节点3"; background_color: "lightpink"}
    ]
}
\stopMindmapExample

\startMindmapExample
\mmjson{
    direction: "top-down";
    name: "折线连接";
    connection: "orthogonal";
    background_color: "lightblue";
    children: [
        {name: "子节点1"; background_color: "lightgreen"},
        {name: "子节点2"; background_color: "lightyellow"},
        {name: "子节点3"; background_color: "lightpink"}
    ]
}
\stopMindmapExample

\startMindmapExample
\mmjson{
    direction: "top-down";
    name: "圆角折线";
    connection: "rounded";
    background_color: "lightblue";
    children: [
        {name: "子节点1"; background_color: "lightgreen"},
        {name: "子节点2"; background_color: "lightyellow"},
        {name: "子节点3"; background_color: "lightpink"}
    ]
}
\stopMindmapExample

\startMindmapExample
\mmjson{
    direction: "left-right";
    name: "横向圆角折线";
    connection: "rounded";
    background_color: "lightblue";
    children: [
        {name: "子节点1"; background_color: "lightgreen"},
        {name: "子节点2"; background_color: "lightyellow"},
        {name: "子节点3"; background_color: "lightpink"}
    ]
}
\stopMindmapExample

\startMindmapExample
\mmjson{
    direction: "bottom-up";
    name: "从下到上圆角";
    connection: "rounded";
    background_color: "lightblue";
    children: [
        {name: "子节点1"; background_color: "lightgreen"},
        {name: "子节点2"; background_color: "lightyellow"},
        {name: "子节点3"; background_color: "lightpink"}
    ]
}
\stopMindmapExample

\startMindmapExample
\mmjson{
    direction: "right-left";
    name: "从右到左圆角";
    connection: "rounded";
    line_radius: 12;
    background_color: "lightblue";
    children: [
        {name: "子节点1"; background_color: "lightgreen"},
        {name: "子节点2"; background_color: "lightyellow"},
        {name: "子节点3"; background_color: "lightpink"}
    ]
}
\stopMindmapExample

\subsubject{4.2 连接线颜色和宽度 (line_color, line_width)}

\startMindmapExample
\mmjson{
    direction: "top-down";
    name: "父节点";
    line_color: "darkblue";
    line_width: 2;
    background_color: "lightblue";
    children: [
        {name: "子节点1"; line_color: "darkgreen"; line_width: 1.5; background_color: "lightgreen"},
        {name: "子节点2"; line_color: "darkorange"; line_width: 1; background_color: "lightyellow"},
        {name: "子节点3"; line_color: "darkred"; line_width: 0.5; background_color: "lightpink"}
    ]
}
\stopMindmapExample

\subsubject{4.3 连接锚点 (connection_anchor)}

\startMindmapExample
\mmjson{
    direction: "top-down";
    name: "边框到边框";
    connection_anchor: "edge-edge";
    background_color: "lightblue";
    children: [
        {name: "子节点1"; background_color: "lightgreen"},
        {name: "子节点2"; background_color: "lightyellow"}
    ]
}
\stopMindmapExample

\startMindmapExample
\mmjson{
    direction: "top-down";
    name: "中心到中心";
    connection_anchor: "center-center";
    background_color: "lightblue";
    children: [
        {name: "子节点1"; background_color: "lightgreen"},
        {name: "子节点2"; background_color: "lightyellow"}
    ]
}
\stopMindmapExample

\subsubject{4.4 锚点位置 (from_pos, to_pos)}

\startframed[frame=off,background=color,backgroundcolor=lightgray,align=flushleft,width=max]
{\bf from\_pos} 和 {\bf to\_pos} 控制连接线在边框上的具体位置。\\
上下结构 (top-down/bottom-up): {\tt left/center/right}\\
左右结构 (left-right/right-left): {\tt top/center/bottom}
\stopframed

\startMindmapExample
\mmjson{
    direction: "top-down";
    name: "左侧连接";
    from_pos: "left";
    to_pos: "left";
    background_color: "lightblue";
    children: [
        {name: "子节点1"; background_color: "lightgreen"},
        {name: "子节点2"; background_color: "lightyellow"}
    ]
}
\stopMindmapExample

\startMindmapExample
\mmjson{
    direction: "top-down";
    name: "右侧连接";
    from_pos: "right";
    to_pos: "right";
    background_color: "lightblue";
    children: [
        {name: "子节点1"; background_color: "lightgreen"},
        {name: "子节点2"; background_color: "lightyellow"}
    ]
}
\stopMindmapExample

\startMindmapExample
\mmjson{
    direction: "left-right";
    name: "顶部连接";
    from_pos: "top";
    to_pos: "top";
    background_color: "lightblue";
    children: [
        {name: "子节点1"; background_color: "lightgreen"},
        {name: "子节点2"; background_color: "lightyellow"}
    ]
}
\stopMindmapExample

\startMindmapExample
\mmjson{
    direction: "left-right";
    name: "底部连接";
    from_pos: "bottom";
    to_pos: "bottom";
    background_color: "lightblue";
    children: [
        {name: "子节点1"; background_color: "lightgreen"},
        {name: "子节点2"; background_color: "lightyellow"}
    ]
}
\stopMindmapExample

\subsubject{4.5 多层节点的曲线连接}

\startMindmapExample
\mmjson{
    direction: "top-down";
    name: "多层曲线";
    connection: "curved";
    background_color: "lightblue";
    children: [
        {
            name: "左分支";
            background_color: "lightgreen";
            children: [
                {name: "左1"; background_color: "white"},
                {name: "左2"; background_color: "white"}
            ]
        },
        {
            name: "中分支";
            background_color: "lightyellow";
            children: [
                {name: "中1"; background_color: "white"},
                {name: "中2"; background_color: "white"}
            ]
        },
        {
            name: "右分支";
            background_color: "lightpink";
            children: [
                {name: "右1"; background_color: "white"},
                {name: "右2"; background_color: "white"}
            ]
        }
    ]
}
\stopMindmapExample

\subsubject{4.6 不同方向的连接类型对比}

\startframed[frame=off,background=color,backgroundcolor=lightgray,align=flushleft,width=max]
{\bf 从上到下 (top-down)} - 直线 / 曲线 / 折线
\stopframed

\startMindmapExample
\mmjson{
    direction: "top-down";
    name: "直线";
    connection: "straight";
    background_color: "lightblue";
    children: [
        {name: "A"; background_color: "lightgreen"},
        {name: "B"; background_color: "lightyellow"},
        {name: "C"; background_color: "lightpink"}
    ]
}
\stopMindmapExample

\startMindmapExample
\mmjson{
    direction: "top-down";
    name: "曲线";
    connection: "curved";
    background_color: "lightblue";
    children: [
        {name: "A"; background_color: "lightgreen"},
        {name: "B"; background_color: "lightyellow"},
        {name: "C"; background_color: "lightpink"}
    ]
}
\stopMindmapExample

\startMindmapExample
\mmjson{
    direction: "top-down";
    name: "折线";
    connection: "orthogonal";
    background_color: "lightblue";
    children: [
        {name: "A"; background_color: "lightgreen"},
        {name: "B"; background_color: "lightyellow"},
        {name: "C"; background_color: "lightpink"}
    ]
}
\stopMindmapExample

\startframed[frame=off,background=color,backgroundcolor=lightgray,align=flushleft,width=max]
{\bf 从下到上 (bottom-up)} - 直线 / 曲线 / 折线
\stopframed

\startMindmapExample
\mmjson{
    direction: "bottom-up";
    name: "直线";
    connection: "straight";
    background_color: "lightblue";
    children: [
        {name: "A"; background_color: "lightgreen"},
        {name: "B"; background_color: "lightyellow"},
        {name: "C"; background_color: "lightpink"}
    ]
}
\stopMindmapExample

\startMindmapExample
\mmjson{
    direction: "bottom-up";
    name: "曲线";
    connection: "curved";
    background_color: "lightblue";
    children: [
        {name: "A"; background_color: "lightgreen"},
        {name: "B"; background_color: "lightyellow"},
        {name: "C"; background_color: "lightpink"}
    ]
}
\stopMindmapExample

\startMindmapExample
\mmjson{
    direction: "bottom-up";
    name: "折线";
    connection: "orthogonal";
    background_color: "lightblue";
    children: [
        {name: "A"; background_color: "lightgreen"},
        {name: "B"; background_color: "lightyellow"},
        {name: "C"; background_color: "lightpink"}
    ]
}
\stopMindmapExample

\startframed[frame=off,background=color,backgroundcolor=lightgray,align=flushleft,width=max]
{\bf 从左到右 (left-right)} - 直线 / 曲线 / 折线
\stopframed

\startMindmapExample
\mmjson{
    direction: "left-right";
    name: "直线";
    connection: "straight";
    background_color: "lightblue";
    children: [
        {name: "A"; background_color: "lightgreen"},
        {name: "B"; background_color: "lightyellow"},
        {name: "C"; background_color: "lightpink"}
    ]
}
\stopMindmapExample

\startMindmapExample
\mmjson{
    direction: "left-right";
    name: "曲线";
    connection: "curved";
    background_color: "lightblue";
    children: [
        {name: "A"; background_color: "lightgreen"},
        {name: "B"; background_color: "lightyellow"},
        {name: "C"; background_color: "lightpink"}
    ]
}
\stopMindmapExample

\startMindmapExample
\mmjson{
    direction: "left-right";
    name: "折线";
    connection: "orthogonal";
    background_color: "lightblue";
    children: [
        {name: "A"; background_color: "lightgreen"},
        {name: "B"; background_color: "lightyellow"},
        {name: "C"; background_color: "lightpink"}
    ]
}
\stopMindmapExample

\startframed[frame=off,background=color,backgroundcolor=lightgray,align=flushleft,width=max]
{\bf 从右到左 (right-left)} - 直线 / 曲线 / 折线
\stopframed

\startMindmapExample
\mmjson{
    direction: "right-left";
    name: "直线";
    connection: "straight";
    background_color: "lightblue";
    children: [
        {name: "A"; background_color: "lightgreen"},
        {name: "B"; background_color: "lightyellow"},
        {name: "C"; background_color: "lightpink"}
    ]
}
\stopMindmapExample

\startMindmapExample
\mmjson{
    direction: "right-left";
    name: "曲线";
    connection: "curved";
    background_color: "lightblue";
    children: [
        {name: "A"; background_color: "lightgreen"},
        {name: "B"; background_color: "lightyellow"},
        {name: "C"; background_color: "lightpink"}
    ]
}
\stopMindmapExample

\startMindmapExample
\mmjson{
    direction: "right-left";
    name: "折线";
    connection: "orthogonal";
    background_color: "lightblue";
    children: [
        {name: "A"; background_color: "lightgreen"},
        {name: "B"; background_color: "lightyellow"},
        {name: "C"; background_color: "lightpink"}
    ]
}
\stopMindmapExample

\subject{五、属性继承}

\subsubject{5.1 基本继承 (inherit)}

\startframed[frame=off,background=color,backgroundcolor=lightgray,align=flushleft,width=max]
{\bf inherit.level} 表示往后继承的层数。设置 inherit 的节点本身不参与计算，从其直接子节点开始往下数 level 层。
\stopframed

\startMindmapExample
\mmjson{
    direction: "left-right";
    name: "根节点";
    background_color: "lightblue";
    border_color: "darkblue";
    inherit: {
        background_color: "lightgreen",
        border_color: "darkgreen",
        font_size: 12,
        level: 2
    };
    children: [
        {
            name: "子节点(第1层)";
            children: [
                {name: "孙节点(第2层)"},
                {name: "曾孙(第3层)"}
            ]
        }
    ]
}
\stopMindmapExample

\subsubject{5.2 从子节点开始继承}

\startMindmapExample
\mmjson{
    direction: "left-right";
    name: "根节点";
    background_color: "lightblue";
    children: [
        {
            name: "子节点";
            background_color: "lightgreen";
            inherit: {
                background_color: "lightyellow",
                border_color: "darkorange",
                level: 1
            };
            children: [
                {
                    name: "孙节点(继承)";
                    children: [
                        {name: "曾孙(不继承)"}
                    ]
                }
            ]
        }
    ]
}
\stopMindmapExample

\subject{六、网格紧凑布局}

\subsubject{6.1 紧凑布局对比}

\startframed[frame=off,background=color,backgroundcolor=lightgray,align=flushleft,width=max]
{\bf enable\_grid\_compact: true} - 启用网格紧凑布局，利用空白空间使布局更紧凑。
\stopframed

\startMindmapExample
\mmjson{
    direction: "left-right";
    enable_grid_compact: true;
    name: "知识体系";
    background_color: "lightblue";
    children: [
        {
            name: "编程语言";
            background_color: "lightgreen";
            children: [
                {name: "Python"; background_color: "white"},
                {name: "JavaScript"; background_color: "white"},
                {name: "PHP"; background_color: "white"}
            ]
        },
        {
            name: "算法数据结构";
            background_color: "lightyellow";
        },
        {
            name: "系统设计";
            background_color: "lightpink";
        }
    ]
}
\stopMindmapExample

\subsubject{6.2 属性兼容性测试}

\startMindmapExample
\mmjson{
    direction: "left-right";
    enable_grid_compact: true;
    name: "属性测试";
    background_color: "lightblue";
    border_color: "darkblue";
    border_width: 2;
    font_size: 12;
    text_color: "darkblue";
    width: 80;
    shape: "rounded";
    corner_radius: 8;
    line_color: "darkblue";
    line_width: 1.5;
    connection: "curved";
    children: [
        {
            name: "样式测试";
            background_color: "lightgreen";
            border_color: "darkgreen";
            border_width: 2;
            shape: "ellipse";
            width: 70;
            children: [
                {name: "圆形"; background_color: "white"; shape: "circle"; width: 50},
                {name: "椭圆"; background_color: "white"; shape: "ellipse"; width: 60}
            ]
        },
        {
            name: "文本测试";
            background_color: "lightyellow";
            font_size: 11;
            text_color: "darkorange";
            text_align: "left";
            width: 70;
        },
        {
            name: "边框测试";
            background_color: "lightpink";
            border_color: "darkred";
            border_width: 3;
            corner_radius: 10;
            width: 70;
        },
        {
            name: "连接线测试";
            background_color: "lightcyan";
            line_color: "darkcyan";
            line_width: 2;
            width: 70;
        }
    ]
}
\stopMindmapExample

\subsubject{6.3 Align 属性测试}

\startMindmapExample
\mmjson{
    direction: "left-right";
    enable_grid_compact: true;
    name: "Align测试";
    background_color: "lightblue";
    width: 70;
    children: [
        {
            name: "Top对齐";
            background_color: "lightgreen";
            align: "top";
            width: 60;
            children: [
                {name: "A1"; background_color: "white"},
                {name: "A2"; background_color: "white"},
                {name: "A3"; background_color: "white"}
            ]
        },
        {
            name: "Center";
            background_color: "lightyellow";
            align: "center";
            width: 60;
        },
        {
            name: "Bottom对齐";
            background_color: "lightpink";
            align: "bottom";
            width: 60;
            children: [
                {name: "B1"; background_color: "white"},
                {name: "B2"; background_color: "white"}
            ]
        }
    ]
}
\stopMindmapExample

\subject{七、综合示例}

\subsubject{7.1 项目管理思维导图}

\startMindmapExample
\mmjson{
    direction: "left-right";
    name: "项目管理";
    background_color: "lightblue";
    border_color: "darkblue";
    font_size: 11;
    width: 80;
    connection: "curved";
    children: [
        {
            name: "计划阶段";
            background_color: "lightgreen";
            width: 80;
            children: [
                {name: "需求分析"; background_color: "white"; width: 70},
                {name: "资源评估"; background_color: "white"; width: 70},
                {name: "时间规划"; background_color: "white"; width: 70}
            ]
        },
        {
            name: "执行阶段";
            background_color: "lightyellow";
            width: 80;
            children: [
                {name: "任务分配"; background_color: "white"; width: 70},
                {name: "进度跟踪"; background_color: "white"; width: 70},
                {name: "风险控制"; background_color: "white"; width: 70}
            ]
        },
        {
            name: "收尾阶段";
            background_color: "lightpink";
            width: 80;
            children: [
                {name: "成果验收"; background_color: "white"; width: 70},
                {name: "文档归档"; background_color: "white"; width: 70},
                {name: "经验总结"; background_color: "white"; width: 70}
            ]
        }
    ]
}
\stopMindmapExample

\subsubject{7.2 软件开发流程}

\startMindmapExample
\mmjson{
    direction: "left-right";
    name: "软件开发";
    background_color: "lightblue";
    border_color: "darkblue";
    font_size: 11;
    width: 80;
    connection: "curved";
    children: [
        {
            name: "需求";
            background_color: "lightgreen";
            width: 70;
            children: [
                {name: "用户需求"; background_color: "white"; width: 60},
                {name: "功能需求"; background_color: "white"; width: 60}
            ]
        },
        {
            name: "设计";
            background_color: "lightyellow";
            width: 70;
            children: [
                {name: "架构设计"; background_color: "white"; width: 60},
                {name: "UI设计"; background_color: "white"; width: 60}
            ]
        },
        {
            name: "开发";
            background_color: "lightpink";
            width: 70;
            children: [
                {name: "前端开发"; background_color: "white"; width: 60},
                {name: "后端开发"; background_color: "white"; width: 60}
            ]
        },
        {
            name: "测试";
            background_color: "lightcyan";
            width: 70;
            children: [
                {name: "单元测试"; background_color: "white"; width: 60},
                {name: "集成测试"; background_color: "white"; width: 60}
            ]
        }
    ]
}
\stopMindmapExample

\subject{七、主题系统}

\subsubject{7.1 定义主题 (definemmtheme)}

\startframed[frame=off,background=color,backgroundcolor=lightgray,align=flushleft,width=max]
{\bf \textbackslash definemmtheme[主题名][属性...]} - 定义可复用的主题样式。
\stopframed

\definemmtheme[blue][
    direction: "left-right",
    connection: "rounded",
    enable_grid_compact: true,
    text_color=darkblue,
    background_color=lightblue,
    border_color=lightblue,
    from_pos=bottom,
    to_pos=bottom,
]

\definemmtheme[green][
    text_color=darkgreen,
    background_color=lightgreen,
    border_color=darkgreen,
]

\definemmtheme[orange][
    text_color=darkorange,
    background_color=lightyellow,
    border_color=darkorange,
]

\startMindmapExample
\mmjson{
    name: "使用主题";
    theme: "blue";
    children: [
        {name: "绿色主题"; theme: "green"},
        {name: "橙色主题"; theme: "orange"},
        {name: "默认样式"}
    ]
}
\stopMindmapExample

\subsubject{7.2 主题覆盖}

\startframed[frame=off,background=color,backgroundcolor=lightgray,align=flushleft,width=max]
主题属性可以被节点级别的属性覆盖。
\stopframed

\startMindmapExample
\mmjson{
    direction: "left-right";
    name: "主题覆盖测试";
    theme: "blue";
    children: [
        {name: "继承主题"},
        {name: "覆盖背景"; background_color: "lightpink"},
        {name: "覆盖字体"; font_size: 14; text_color: "darkred"},
        {name: "覆盖形状"; shape: "ellipse"; background_color: "lightcyan"}
    ]
}
\stopMindmapExample

\subsubject{7.3 修改主题 (setupmmtheme)}

\startframed[frame=off,background=color,backgroundcolor=lightgray,align=flushleft,width=max]
{\bf \textbackslash setupmmtheme[主题名][属性...]} - 修改已定义的主题。
\stopframed

\definemmtheme[custom][
    text_color=black,
    background_color=white,
    border_color=gray,
]

\setupmmtheme[custom][
    text_color=darkblue,
    background_color=lightblue,
    shape=rounded,
    corner_radius=5,
]

\startMindmapExample
\mmjson{
    direction: "top-down";
    name: "修改后的主题";
    theme: "custom";
    children: [
        {name: "子节点1"},
        {name: "子节点2"},
        {name: "子节点3"}
    ]
}
\stopMindmapExample

\subsubject{7.4 主题继承链}

\startMindmapExample
\mmjson{
    direction: "left-right";
    name: "根节点";
    theme: "blue";
    children: [
        {
            name: "分支1";
            theme: "green";
            children: [
                {name: "孙节点1"},
                {name: "孙节点2"}
            ]
        },
        {
            name: "分支2";
            theme: "orange";
            children: [
                {name: "孙节点3"},
                {name: "孙节点4"}
            ]
        }
    ]
}
\stopMindmapExample

\subsubject{7.5 行高设置 (lineheight)}

\definemmtheme[compact][
    lineheight=1.0,
]

\definemmtheme[spacious][
    lineheight=2.0,
]

\startMindmapExample
\mmjson{
    direction: "top-down";
    name: "紧凑行高";
    theme: "compact";
    width: 80;
    background_color: "lightblue";
    children: [
        {name: "这是一个较长的文本用于测试行高效果"; width: 70; background_color: "lightgreen"},
        {name: "另一段长文本用于测试多行显示效果"; width: 70; background_color: "lightyellow"}
    ]
}
\stopMindmapExample

\startMindmapExample
\mmjson{
    direction: "top-down";
    name: "宽松行高";
    theme: "spacious";
    width: 80;
    background_color: "lightblue";
    children: [
        {name: "这是一个较长的文本用于测试行高效果"; width: 70; background_color: "lightgreen"},
        {name: "另一段长文本用于测试多行显示效果"; width: 70; background_color: "lightyellow"}
    ]
}
\stopMindmapExample

\subsubject{7.6 主题继承 (definemmtheme 继承)}

\startframed[frame=off,background=color,backgroundcolor=lightgray,align=flushleft,width=max]
{\bf \textbackslash definemmtheme[新主题][父主题]} - 新主题继承父主题的所有属性。
\stopframed

\definemmtheme[basedark][
    text_color=white,
    background_color=darkgray,
    border_color=black,
    shape=rounded,
    corner_radius=8,
]

\definemmtheme[lightblue][basedark]
\setupmmtheme[lightblue][
    background_color=lightblue,
    text_color=darkblue,
    border_color=blue,
]

\definemmtheme[lightgreen][basedark]
\setupmmtheme[lightgreen][
    background_color=lightgreen,
    text_color=darkgreen,
    border_color=green,
]

\startMindmapExample
\mmjson{
    direction: "left-right";
    name: "继承测试";
    theme: "basedark";
    children: [
        {name: "浅蓝主题"; theme: "lightblue"},
        {name: "浅绿主题"; theme: "lightgreen"},
        {name: "原始深色"}
    ]
}
\stopMindmapExample

\subsubject{7.7 层级主题 (definemmtheme 三参数)}

\startframed[frame=off,background=color,backgroundcolor=lightgray,align=flushleft,width=max]
{\bf \textbackslash definemmtheme[主题名][层级][属性...]} - 为不同层级定义不同样式。\\
层级 0 = 根节点，层级 1 = 第一层子节点，以此类推。
\stopframed

\definemmtheme[leveltheme]
\definemmtheme[leveltheme][0][
    text_color=white,
    background_color=darkblue,
    shape=ellipse,
]
\definemmtheme[leveltheme][1][
    text_color=darkblue,
    background_color=lightblue,
    shape=rounded,
    corner_radius=10,
]
\definemmtheme[leveltheme][2][
    text_color=darkgreen,
    background_color=lightgreen,
    shape=rectangle,
]
\definemmtheme[leveltheme][3][
    text_color=darkorange,
    background_color=lightyellow,
    shape=rounded,
]

\startMindmapExample
\mmjson{
    direction: "top-down";
    name: "层级主题根节点";
    theme: "leveltheme";
    children: [
        {
            name: "第一层";
            children: [
                {
                    name: "第二层";
                    children: [
                        {name: "第三层A"},
                        {name: "第三层B"}
                    ]
                }
            ]
        },
        {
            name: "第一层B";
            children: [
                {name: "第二层C"},
                {name: "第二层D"}
            ]
        }
    ]
}
\stopMindmapExample

\subsubject{7.8 层级主题继承}

\startframed[frame=off,background=color,backgroundcolor=lightgray,align=flushleft,width=max]
结合主题继承和层级设置，创建更灵活的主题系统。
\stopframed

\definemmtheme[baselevel][
    direction: "left-right",
    connection: "rounded",
]
\definemmtheme[baselevel][0][
    text_color=white,
    background_color=purple,
    shape=ellipse,
]
\definemmtheme[baselevel][1][
    background_color=plum,
    text_color=purple,
]

\definemmtheme[childlevel][baselevel]
\setupmmtheme[childlevel][0][
    background_color=darkred,
]
\setupmmtheme[childlevel][1][
    background_color=lightcoral,
    text_color=darkred,
]

\startMindmapExample
\mmjson{
    name: "父主题";
    theme: "baselevel";
    children: [
        {name: "子节点1"},
        {name: "子节点2"}
    ]
}
\stopMindmapExample

\startMindmapExample
\mmjson{
    name: "子主题";
    theme: "childlevel";
    children: [
        {name: "子节点1"},
        {name: "子节点2"}
    ]
}
\stopMindmapExample

\subject{六、思维导图布局}

\subsubject{6.1 中心节点左右均衡分布}

\startframed[frame=off,background=color,backgroundcolor=lightgray,align=flushleft,width=max]
{\bf direction: "mindmap"} 将子节点均匀分布在中心节点的左右两侧，奇数子节点在左侧，偶数子节点在右侧。
\stopframed

\startMindmapExample
\mmjson{
    direction: "mindmap";
    name: "中心主题";
    background_color: "lightblue";
    border_color: "darkblue";
    children: [
        {name: "分支1"; background_color: "lightgreen"; children: [{name: "子1"}, {name: "子2"}]},
        {name: "分支2"; background_color: "lightyellow"; children: [{name: "子1"}, {name: "子2"}]},
        {name: "分支3"; background_color: "lightpink"; children: [{name: "子1"}]},
        {name: "分支4"; background_color: "lightcyan"; children: [{name: "子1"}, {name: "子2"}]}
    ]
}
\stopMindmapExample

\subsubject{6.2 思维导图布局的曲线连接}

\startMindmapExample
\mmjson{
    direction: "mindmap";
    connection: "curved";
    name: "项目";
    background_color: "lightblue";
    children: [
        {name: "需求"; background_color: "lightgreen"},
        {name: "设计"; background_color: "lightyellow"},
        {name: "开发"; background_color: "lightpink"},
        {name: "测试"; background_color: "lightcyan"}
    ]
}
\stopMindmapExample

\subject{七、连线端点装饰}

\subsubject{7.1 箭头端点}

\startframed[frame=off,background=color,backgroundcolor=lightgray,align=flushleft,width=max]
{\bf line\_start\_marker} 和 {\bf line\_end\_marker} 可设置连线端点的装饰类型。\\
可选值: {\tt arrow} (箭头), {\tt circle/dot} (圆点), {\tt square} (方点), {\tt diamond} (菱形)
\stopframed

\startMindmapExample
\mmjson{
    direction: "left-right";
    name: "流程";
    background_color: "lightblue";
    line_end_marker: "arrow";
    children: [
        {name: "步骤1"; background_color: "lightgreen"; line_end_marker: "arrow"},
        {name: "步骤2"; background_color: "lightyellow"; line_end_marker: "arrow"},
        {name: "步骤3"; background_color: "lightpink"}
    ]
}
\stopMindmapExample

\subsubject{7.2 圆点端点}

\startMindmapExample
\mmjson{
    direction: "top-down";
    name: "任务列表";
    background_color: "lightblue";
    line_start_marker: "circle";
    line_end_marker: "circle";
    children: [
        {name: "任务A"; background_color: "lightgreen"},
        {name: "任务B"; background_color: "lightyellow"},
        {name: "任务C"; background_color: "lightpink"}
    ]
}
\stopMindmapExample

\subsubject{7.3 方点端点}

\startMindmapExample
\mmjson{
    direction: "left-right";
    name: "模块";
    background_color: "lightblue";
    line_end_marker: "square";
    children: [
        {name: "模块A"; background_color: "lightgreen"},
        {name: "模块B"; background_color: "lightyellow"}
    ]
}
\stopMindmapExample

\subsubject{7.4 菱形端点}

\startMindmapExample
\mmjson{
    direction: "top-down";
    name: "决策";
    background_color: "lightblue";
    line_end_marker: "diamond";
    children: [
        {name: "选项1"; background_color: "lightgreen"},
        {name: "选项2"; background_color: "lightyellow"}
    ]
}
\stopMindmapExample

\subsubject{7.5 混合端点}

\startMindmapExample
\mmjson{
    direction: "left-right";
    name: "数据流";
    background_color: "lightblue";
    line_start_marker: "circle";
    line_end_marker: "arrow";
    children: [
        {name: "处理1"; background_color: "lightgreen"; line_start_marker: "circle"; line_end_marker: "arrow"},
        {name: "处理2"; background_color: "lightyellow"; line_start_marker: "circle"; line_end_marker: "arrow"},
        {name: "输出"; background_color: "lightpink"; line_start_marker: "circle"}
    ]
}
\stopMindmapExample

\subject{八、括号连接}

\subsubject{8.1 左括号连接}

\startframed[frame=off,background=color,backgroundcolor=lightgray,align=flushleft,width=max]
{\bf connection: "brace-left"} 使用左括号形状 {\tt \{} 包围所有子节点。\\
括号从父节点延伸，覆盖所有子节点的垂直范围，不单独连接每个子节点。
\stopframed

\startMindmapExample
\mmjson{
    direction: "left-right";
    connection: "brace-left";
    name: "分类";
    background_color: "lightblue";
    children: [
        {name: "项目A"; background_color: "lightgreen"},
        {name: "项目B"; background_color: "lightyellow"},
        {name: "项目C"; background_color: "lightpink"}
    ]
}
\stopMindmapExample

\subsubject{8.2 右括号连接}

\startMindmapExample
\mmjson{
    direction: "left-right";
    connection: "brace-right";
    name: "分类";
    background_color: "lightblue";
    children: [
        {name: "项目A"; background_color: "lightgreen"},
        {name: "项目B"; background_color: "lightyellow"},
        {name: "项目C"; background_color: "lightpink"}
    ]
}
\stopMindmapExample

\subject{九、节点偏移}

\subsubject{9.1 X 轴偏移 (offset_x)}

% \startframed[frame=off,background=color,backgroundcolor=lightgray,align=flushleft,width=max]
% {\bf offset\_x} 控制节点在 X 轴方向的偏移量（单位：pt）。\\
% 正值向右偏移，负值向左偏移，默认值为 0。
% \stopframed

\startMindmapExample
\mmjson{
    direction: "top-down";
    name: "中心主题";
    background_color: "lightblue";
    children: [
        {
            name: "向右偏移";
            background_color: "lightgreen";
            offset_x: 50;
            children: [
                {name: "子节点1"; background_color: "white"},
                {name: "子节点2"; background_color: "white"}
            ]
        },
        {
            name: "向左偏移";
            background_color: "lightyellow";
            offset_x: -30;
            children: [
                {name: "子节点1"; background_color: "white"},
                {name: "子节点2"; background_color: "white"}
            ]
        },
        {
            name: "无偏移";
            background_color: "lightpink";
            children: [
                {name: "子节点1"; background_color: "white"},
                {name: "子节点2"; background_color: "white"}
            ]
        }
    ]
}
\stopMindmapExample

\subsubject{9.2 Y 轴偏移 (offset_y)}

% \startframed[frame=off,background=color,backgroundcolor=lightgray,align=flushleft,width=max]
% {\bf offset\_y} 控制节点在 Y 轴方向的偏移量（单位：pt）。\\
% 正值向下偏移，负值向上偏移，默认值为 0。
% \stopframed

\startMindmapExample
\mmjson{
    direction: "top-down";
    name: "中心主题";
    background_color: "lightblue";
    children: [
        {
            name: "向下偏移";
            background_color: "lightgreen";
            offset_y: 40;
            children: [
                {name: "子节点1"; background_color: "white"},
                {name: "子节点2"; background_color: "white"}
            ]
        },
        {
            name: "向上偏移";
            background_color: "lightyellow";
            offset_y: -20;
            children: [
                {name: "子节点1"; background_color: "white"},
                {name: "子节点2"; background_color: "white"}
            ]
        }
    ]
}
\stopMindmapExample

\subsubject{9.3 同时偏移 X 轴和 Y 轴}

\startMindmapExample
\mmjson{
    direction: "top-down";
    name: "中心主题";
    background_color: "lightblue";
    children: [
        {
            name: "右下偏移";
            background_color: "lightgreen";
            offset_x: 40;
            offset_y: 30;
            children: [
                {name: "子节点1"; background_color: "white"}
            ]
        },
        {
            name: "左上偏移";
            background_color: "lightyellow";
            offset_x: -50;
            offset_y: -25;
            children: [
                {name: "子节点1"; background_color: "white"}
            ]
        }
    ]
}
\stopMindmapExample

\subsubject{9.4 深层节点偏移}

\startframed[frame=off,background=color,backgroundcolor=lightgray,align=flushleft,width=max]
偏移功能支持任意深度的节点，可以为不同层级的节点设置不同的偏移量。
\stopframed

\startMindmapExample
\mmjson{
    direction: "top-down";
    name: "根节点";
    background_color: "lightblue";
    children: [
        {
            name: "第一层";
            background_color: "lightgreen";
            offset_x: 30;
            children: [
                {
                    name: "第二层";
                    background_color: "lightyellow";
                    offset_y: 20;
                    children: [
                        {
                            name: "第三层";
                            background_color: "lightpink";
                            offset_x: -20;
                            children: [
                                {name: "第四层"; background_color: "white"}
                            ]
                        }
                    ]
                }
            ]
        }
    ]
}
\stopMindmapExample

\subsubject{9.5 横向布局偏移}

\startMindmapExample
\mmjson{
    direction: "left-right";
    name: "中心主题";
    background_color: "lightblue";
    children: [
        {
            name: "X 轴偏移";
            background_color: "lightgreen";
            offset_x: 50;
            children: [
                {name: "子节点1"; background_color: "white"}
            ]
        },
        {
            name: "Y 轴偏移";
            background_color: "lightyellow";
            offset_y: -30;
            children: [
                {name: "子节点1"; background_color: "white"}
            ]
        }
    ]
}
\stopMindmapExample

\subsubject{9.7 微调节点位置}

\startframed[frame=off,background=color,backgroundcolor=lightgray,align=flushleft,width=max]
使用小偏移量进行微调，可以优化节点的视觉布局，避免重叠或改善美观度。
\stopframed

\startMindmapExample
\mmjson{
    direction: "top-down";
    name: "项目计划";
    background_color: "lightblue";
    children: [
        {
            name: "需求分析";
            background_color: "lightgreen";
            offset_x: 10;
            children: [
                {name: "用户调研"; background_color: "white"},
                {name: "需求文档"; background_color: "white"}
            ]
        },
        {
            name: "设计阶段";
            background_color: "lightyellow";
            offset_x: -5;
            children: [
                {name: "UI设计"; background_color: "white"},
                {name: "架构设计"; background_color: "white"}
            ]
        },
        {
            name: "开发实现";
            background_color: "lightpink";
            offset_y: 5;
            children: [
                {name: "前端开发"; background_color: "white"},
                {name: "后端开发"; background_color: "white"}
            ]
        }
    ]
}
\stopMindmapExample

\subsubject{9.8 不同布局算法的偏移}

\startframed[frame=off,background=color,backgroundcolor=lightgray,align=flushleft,width=max]
偏移功能支持所有布局算法：hybrid、scanline、dp、radial。
\stopframed

\startMindmapExample
\mmjson{
    direction: "top-down";
    method: "hybrid";
    name: "混合布局";
    background_color: "lightblue";
    children: [
        {name: "节点1"; background_color: "lightgreen"; offset_x: 30},
        {name: "节点2"; background_color: "lightyellow"; offset_y: 20}
    ]
}
\stopMindmapExample

\startMindmapExample
\mmjson{
    direction: "top-down";
    method: "scanline";
    name: "扫描线布局";
    background_color: "lightblue";
    children: [
        {name: "节点1"; background_color: "lightgreen"; offset_x: 30},
        {name: "节点2"; background_color: "lightyellow"; offset_y: 20}
    ]
}
\stopMindmapExample

\startMindmapExample
\mmjson{
    direction: "top-down";
    method: "dp";
    name: "动态规划布局";
    background_color: "lightblue";
    children: [
        {name: "节点1"; background_color: "lightgreen"; offset_x: 30},
        {name: "节点2"; background_color: "lightyellow"; offset_y: 20}
    ]
}
\stopMindmapExample

\subsubject{9.9 偏移量限制}

\startframed[frame=off,background=color,backgroundcolor=lightgray,align=flushleft,width=max]
单个偏移量的绝对值不能超过 1000pt。超过限制的值会被自动截断并显示警告信息。
\stopframed

\startMindmapExample
\mmjson{
    direction: "top-down";
    name: "中心主题";
    background_color: "lightblue";
    children: [
        {
            name: "大偏移";
            background_color: "lightgreen";
            offset_x: 500;
            children: [
                {name: "子节点"; background_color: "white"}
            ]
        },
        {
            name: "负大偏移";
            background_color: "lightyellow";
            offset_y: -400;
            children: [
                {name: "子节点"; background_color: "white"}
            ]
        }
    ]
}
\stopMindmapExample

\subsubject{9.10 偏移与连接线}

\startframed[frame=off,background=color,backgroundcolor=lightgray,align=flushleft,width=max]
偏移不会自动调整连接线。如果偏移量过大，可能会导致连接线看起来不自然。建议使用较小的偏移量进行微调。
\stopframed

\startMindmapExample
\mmjson{
    direction: "top-down";
    connection: "curved";
    name: "中心主题";
    background_color: "lightblue";
    children: [
        {
            name: "适度偏移";
            background_color: "lightgreen";
            offset_x: 20;
            offset_y: 10;
            children: [
                {name: "子节点1"; background_color: "white"},
                {name: "子节点2"; background_color: "white"}
            ]
        },
        {
            name: "适度偏移";
            background_color: "lightyellow";
            offset_x: -20;
            offset_y: -10;
            children: [
                {name: "子节点1"; background_color: "white"},
                {name: "子节点2"; background_color: "white"}
            ]
        }
    ]
}
\stopMindmapExample

\stoptext
